Large language models (LLMs) are increasingly used in multi-turn settings in
which a task unfolds over many exchanges. Their own earlier replies and an
accumulating task state become part of the context they reason from. We ask
whether this degrades reasoning success, and whether the degradation can be
detected early from measures of reasoning quality.

We combine two complementary, prospectively frozen studies. Study~1 uses short,
sharded math tasks from the Lost in Conversation benchmark. Each of three models
answers every turn twice with identical user information, with and without its
own prior replies. Study~2 is a purpose-built, 100-turn benchmark with exact
per-turn answers, calibrated difficulty and scheduled difficulty shifts. Its
dialogues are answered either from the full dialogue history or from a verified
state summary.

Retaining assistant replies lowered final success in Study~1, significantly for
two of three models. In Study~2, accuracy under full history fell sharply within
about twenty turns, even when difficulty was held fixed. A verified state
summary preserved accuracy in all three models. Neither the
reasoning-quality metrics we examined nor a paired process measure warned of
these failures in a robust, model-general way. Coherence and stability even
rose as accuracy fell.

Long interactions can therefore degrade LLM reasoning in ways that single-stream
quality monitoring does not reveal. Compact, verified state representations are
a promising mitigation.
